\subsubsection*{Lemma 5.3.}
\underline{Let $X$ be a smooth connected quasi-projective variety of
dimension $n$, and let $M$ be a locally free sheaf of finite rank $a$
on $X$.  Let $\xi$ denote the invertible sheaf
$\mathcal O_{\mathbb P(M)}(1)$ on $\mathbb P(M)$, and let
$\pi:\mathbb P(M)\to X$ be the structural morphism.  Recall that
$\operatorname{CH}(\mathbb P(M))$ is a free
$\operatorname{CH}(X)$-module (via $\pi^*$), with basis
$1,\xi,\ldots,\xi^{a-1}$, where $\xi$ satisfies}
\begin{equation}\tag{5.3.1}
\xi^a-C_1\xi^{a-1}+C_2\xi^{a-2}-C_3\xi^{a-3}+\cdots=0,
\end{equation}
\underline{the $C_i$ being the Chern classes of $M$ in
$\operatorname{CH}(X)$.  Put
$C(M)=1+C_1+C_2+\cdots$, the total Chern class of $M$; it is an
invertible element of $\operatorname{CH}(X)$.  Then:}

\begin{enumerate}
\item[(i)] \underline{one has}
\begin{equation}\tag{5.3.2}
\pi_*\!\left(\frac{\xi^{a-1}}{1+\xi}\right)
=\frac{1}{C(M)};
\end{equation}
\item[(ii)] \underline{let
$F(T)\in\operatorname{CH}(X)[T]$ be a polynomial such that the
coefficient of $T^i$ is homogeneous of degree $n+a-1-i$.  Thus
$F(\xi)$ has maximal degree
$n+a-1=\dim\mathbb P(M)$ in
$\operatorname{CH}(\mathbb P(M))$.  Then}
\begin{equation}\tag{5.3.3}
\int_{\mathbb P(M)}F(\xi)
=(-1)^{a-1}\int_X\frac{F(-1)}{C(M)}.
\end{equation}
\end{enumerate}

\noindent\underline{Proof.}
We first deduce (5.3.3) from (5.3.2).  Write
\begin{equation}\tag{5.3.4}
\frac{1}{C(M)}=\sum_{j\geq0}b_j,\qquad
b_0=1,\qquad
b_j\in\operatorname{CH}^j(X),\qquad
b_j=0\ \text{for }j<0.
\end{equation}
Then (5.3.2) is equivalent to
\begin{equation}\tag{5.3.5}
\pi_*(\xi^{a-1+j})=(-1)^jb_j
\qquad\text{whenever }a-1+j\geq0,
\end{equation}
since $\pi_*(\xi^q)=0$ for $q<a-1$ for dimensional reasons.

By linearity it suffices to prove (5.3.3) when
$F(T)=uT^{a-1+j}$, with $u\in\operatorname{CH}^{n-j}(X)$
homogeneous.  By (5.3.5) and the projection formula,
\begin{equation}\tag{5.3.6}
\begin{aligned}
\int_{\mathbb P(M)}F(\xi)
&=\int_{\mathbb P(M)}\pi^*(u)\xi^{a-1+j}\\
&=\int_Xu\,\pi_*(\xi^{a-1+j})
=(-1)^j\int_Xu\,b_j.
\end{aligned}
\end{equation}
On the other hand,
\begin{equation}\tag{5.3.7}
\begin{aligned}
(-1)^{a-1}\int_X\frac{F(-1)}{C(M)}
&=(-1)^{a-1}\int_X
  \frac{(-1)^{a-1+j}u}{C(M)}\\
&=(-1)^j\int_X\sum_i u\,b_i.
\end{aligned}
\end{equation}
Because $\deg u=n-j$ and $\deg b_i=i$, one has
\[
\int_Xu\,b_i=0\qquad\text{for }i\ne j.
\]
Thus (5.3.7) reduces to
\begin{equation}\tag{5.3.8}
(-1)^{a-1}\int_X\frac{F(-1)}{C(M)}
=(-1)^j\int_Xu\,b_j,
\end{equation}
which proves the asserted implication.

It remains to prove (5.3.2), or equivalently (5.3.5).

